\documentclass[a4paper,12pt]{article}

\usepackage[utf8]{inputenc} 
\usepackage{amsmath}
\usepackage[a4paper, margin=0.9in]{geometry} % Adjust the margin here
\usepackage{multirow}
\usepackage{graphicx}
    

\begin{document}

\maketitle
\vspace{-1.5em}


\section*{Question 1}
Consider a system of three point particles, each having an identical mass $m$, with their positions defined by time-dependent vectors,
\begin{equation*}
    \vec r_{1}~=~(at^2,~vt,~0),~\vec r_{2}~=~(vt,~at^2,~b),~\vec r_{3}~=~(0,~b,~at^2), 
\end{equation*}
where $a$, $b$, and $v$ are positive constants. Determine the total external force and the total torque exerted on this system concerning the origin of coordinates. Additionally, based on these findings, analyze whether the system conserves linear and angular momentum.

\section*{Question 2}
Two boats, each having a total mass $M$, are moving in the same direction at a uniform velocity $v$. At a certain point, a package of mass $m$ is thrown from the front boat to the rear boat, with a velocity $u$ \underline{relative to the front boat}. Determine the resulting velocities of both boats and their relative velocity after the package is thrown. Additionally, analyze whether this inelastic process might lead to a collision between the boats.

\section*{Question 3}
Determine the moment of inertia for a thin non-homogeneous rod of length $l$. The rod's linear mass density exhibits a linear increase from $\lambda_1$ to $\lambda_2$ along its length. Calculate the moment of inertia concerning an axis passing through the rod's center and perpendicular to its length.

\section*{Question 4}
Calculate the acceleration of a uniform loop as it \underline{rolls without sliding} on a rough inclined plane inclined at an angle $\alpha$ to the horizontal, in a uniform gravitational field of magnitude $g$.
\end{document}
